\ifdefined\gkver\else\def\gkver{student}\fi
\documentclass[\gkver]{gaokaozhenti}
\begin{document}

\chapter{2018年江苏}

\begin{enumerate}
\item
%% number: 1
%% paperName: 2018年江苏高考第 1 题 3 分
%% typeId: 1
%% type: 单选题
%% chapter: 曲线运动
%% point: 人造地球卫星
%% method: 无
%% score: 3
%% degree: 750
%% duplicateId: 0
%% body:
我国高分系列卫星的高分辨对地观察能力不断提高。今年 5 月 9 日发射的“高分五号”轨道高度约为 $705\Ukm$，之前已运行的“高分四号”轨道高度约为 $36\,000\Ukm$，它们都绕地球做圆周运动。与“高分四号”相比，下列物理量中“高分五号”较小的是\xzanswer{A}

\fourchoices[answer=A]
{周期}
{角速度}
{线速度}
{向心加速度}

%% answer: A
%% memo:
\memoanswer{
设地球质量为 $M$，人造卫星质量为 $m$，人造卫星做匀速圆周运动时，由万有引力提供向心力得 $G\frac{Mm}{r^{2}}=m\frac{v^{2}}{r}=m\omega^{2}r=m\left(\frac{2\pi}{T}\right)^{2}r=ma$，解得 $v=\sqrt{\frac{GM}{r}}$，$\omega=\sqrt{\frac{GM}{r^{3}}}$，$T=2\pi\sqrt{\frac{r^{3}}{GM}}$，$a=\frac{GM}{r^{2}}$，因“高分四号”的轨道半径比“高分五号”的轨道半径大，故 A 正确，B、C、D 错误。
}

\item
%% number: 2
%% paperName: 2018年江苏高考第 2 题 3 分
%% typeId: 1
%% type: 单选题
%% chapter: 电磁感应
%% point: 远距离输电
%% method: 无
%% score: 3
%% degree: 750
%% duplicateId: 0
%% body:
采用 $220\UkV$ 高压向远方的城市输电。当输送功率一定时，为使输电线上损耗的功率减小为原来的 $\frac{1}{4}$，输电电压应变为\xzanswer{C}

\fourchoices[answer=C]
{$55\UkV$}
{$110\UkV$}
{$440\UkV$}
{$880\UkV$}

%% answer: C
%% memo:
\memoanswer{
输电功率 $P=UI$，$U$ 为输电电压，$I$ 为输电线路中的电流，输电线路损失的功率为 $P_{\text{损}}=I^{2}R$，$R$ 为输电线路的电阻，即 $P_{\text{损}}=\left(\frac{P}{U}\right)^{2}R$。当输电功率一定时，输电线路损失的功率为原来的 $\frac{1}{4}$，则输电电压为原来的 2 倍，即 $440\UkV$，故 C 正确。
}

\item
%% number: 3
%% paperName: 2018年江苏高考第 3 题 3 分
%% typeId: 1
%% type: 单选题
%% chapter: 曲线运动
%% point: 平抛运动
%% method: 无
%% score: 3
%% degree: 750
%% duplicateId: 0
%% body:
某弹射管每次弹出的小球速度相等。在沿光滑竖直轨道自由下落过程中，该弹射管保持水平，先后弹出两只小球。忽略空气阻力，两只小球落到水平地面的\xzanswer{B}

\fourchoices[answer=B]
{时刻相同，地点相同}
{时刻相同，地点不同}
{时刻不同，地点相同}
{时刻不同，地点不同}

%% answer: B
%% memo:
\memoanswer{
小球弹出后做曲线运动，该运动可分解成水平方向的匀速运动和竖直方向的匀加速直线运动，竖直方向两球速度相同并且始终在同一高度上，水平方向速度也相同，由于后弹出的小球在空中运动的时间较短，所以两球落地的时刻相同但水平方向通过的距离不同，所以落地点不同，故 B 正确。
}

\item
%% number: 4
%% paperName: 2018年江苏高考第 4 题 3 分
%% typeId: 1
%% type: 单选题
%% chapter: 直线运动
%% point: 竖直上下抛运动
%% method: 无
%% score: 3
%% degree: 550
%% duplicateId: 0
%% body:
从地面竖直向上抛出一只小球，小球运动一段时间后落回地面。忽略空气阻力，该过程中小球的动能 $E_{k}$ 与时间 $t$ 的关系图像是\xzanswer{A}

\fourchoices[answer=A, ispicture=true, h=2.4cm]
{figs/04a.png}
{figs/04b.png}
{figs/04c.png}
{figs/04d.png}

%% answer: A
%% memo:
\memoanswer{
小球做竖直上抛运动时，速度 $v=v_{0}-gt$，由动能 $E_{k}=\frac{1}{2}mv^{2}$ 得 $E_{k}=\frac{1}{2}m(v_{0}-gt)^{2}$，故图象 A 正确。
}

\item
%% number: 5
%% paperName: 2018年江苏高考第 5 题 3 分
%% typeId: 1
%% type: 单选题
%% chapter: 电场
%% point: 带电粒子的平衡
%% method: 无
%% score: 3
%% degree: 450
%% duplicateId: 0
%% body:
如图所示，水平金属板 A、B 分别与电源两极相连，带电油滴处于静止状态。现将 B 板右端向下移动一小段距离，两金属板表面仍均为等势面，则该油滴\xzanswer{D}

\onepicture[w=5.0cm]{figs/05.png}

\fourchoices[answer=D]
{仍然保持静止}
{竖直向下运动}
{向左下方运动}
{向右下方运动}

%% answer: D
%% memo:
\memoanswer{
两极板平行时带电粒子处于平衡状态，则重力等于电场力；当下极板旋转时，板间距离增大，场强减小，电场力小于重力；由于电场线垂直于金属板表面，所以电荷处的电场线如图所示，所以重力与电场力的合力偏向右下方，故粒子向右下方运动，故 D 正确。

\onepicture[w=4.0cm]{figs/05_an.jpg}
}

\item
%% number: 6
%% paperName: 2018年江苏高考第 6 题 4 分
%% typeId: 2
%% type: 多选题
%% chapter: 曲线运动
%% point: 线速度角速度
%% method: 无
%% score: 4
%% degree: 650
%% duplicateId: 0
%% body:
火车以 $60\Ums$ 的速率转过一段弯道，某乘客发现放在桌面上的指南针在 $10\Us$ 内匀速转过了约 $10^{\circ}$。在此 $10\Us$ 时间内，火车\xzanswer{AD}

\fourchoices[answer=AD]
{运动路程为 $600\Um$}
{加速度为零}
{角速度约为 $1\Urads$}
{转弯半径约为 $3.4\Ukm$}

%% answer: AD
%% memo:
\memoanswer{
圆周运动的弧长 $s=vt=60\Ums\times10\Us=600\Um$，故 A 正确；火车转弯是圆周运动，圆周运动是变速运动，所以合力不为零，加速度不为零，故 B 错误；由题意得圆周运动的角速度 $\omega=\frac{\Delta\theta}{\Delta t}=\frac{\frac{10^{\circ}}{360^{\circ}}\times2\pi}{10\Us}=\frac{3.14}{180}\Urads$，又 $v=\omega r$，所以 $r=\frac{v}{\omega}=\frac{60}{3.14}\times180\Um=3439\Um$，故 C 错误，D 正确。
}

\item
%% number: 7
%% paperName: 2018年江苏高考第 7 题 4 分
%% typeId: 2
%% type: 多选题
%% chapter: 运动和力
%% point: 变加速直线运动
%% method: 无
%% score: 4
%% degree: 450
%% duplicateId: 0
%% body:
如图所示，轻质弹簧一端固定，另一端连接一小物块，O 点为弹簧在原长时物块的位置。物块由 A 点静止释放，沿粗糙程度相同的水平面向右运动，最远到达 B 点。在从 A 到 B 的过程中，物块\xzanswer{AD}

\onepicture[w=6.0cm]{figs/07.png}

\fourchoices[answer=AD]
{加速度先减小后增大}
{经过 O 点时的速度最大}
{所受弹簧弹力始终做正功}
{所受弹簧弹力做的功等于克服摩擦力做的功}

%% answer: AD
%% memo:
\memoanswer{
物体从 A 点到 O 点过程，弹力逐渐减为零，刚开始弹簧弹力大于摩擦力，故可分为弹力大于摩擦力过程和弹力小于摩擦力过程：弹力大于摩擦力过程，合力向右，加速度也向右，由于弹力减小，摩擦力不变，小球所受合力减小，加速度减小，弹力等于摩擦力时速度最大，此位置在 A 点和 O 点之间；弹力小于摩擦力过程，合力方向与运动方向相反，弹力减小，摩擦力大小不变，物体所受合力增大，物体的加速度随弹簧形变量的减小而增加，物体做减速运动；从 O 点到 B 点的过程弹力增大，合力向左，加速度继续增大，故 A 正确，B 错误；从 A 点到 O 点过程，弹簧由压缩恢复原长，弹力做正功，从 O 点到 B 点的过程，弹簧伸长，弹力做负功，故 C 错误；从 A 到 B 的过程中由动能定理知，弹簧弹力做的功等于物体克服摩擦力做的功，故 D 正确。
}

\item
%% number: 8
%% paperName: 2018年江苏高考第 8 题 4 分
%% typeId: 2
%% type: 多选题
%% chapter: 电路
%% point: 电容
%% method: 无
%% score: 4
%% degree: 550
%% duplicateId: 0
%% body:
如图所示，电源 $E$ 对电容器 $C$ 充电，当 $C$ 两端电压达到 $80\UV$ 时，闪光灯瞬间导通并发光，$C$ 放电。放电后，闪光灯断开并熄灭，电源再次对 $C$ 充电。这样不断地充电和放电，闪光灯就周期性地发光。该电路\xzanswer{BCD}

\onepicture[w=4.6cm]{figs/08.png}

\fourchoices[answer=BCD]
{充电时，通过 $R$ 的电流不变}
{若 $R$ 增大，则充电时间变长}
{若 $C$ 增大，则闪光灯闪光一次通过的电荷量增大}
{若 $E$ 减小为 $85\UV$，闪光灯闪光一次通过的电荷量不变}

%% answer: BCD
%% memo:
\memoanswer{
电容器充电时两端电压不断增大，所以电源与电容器极板间的电势差不断减小，因此充电电流变小，故 A 错误；当电阻 $R$ 增大时，充电电流变小，电容器所充电荷量不变的情况下，充电时间变长，故 B 正确；若 $C$ 增大，由 $Q=CU$，电容器的带电荷量增大，故 C 正确；当电源电动势为 $85\UV$ 时，电源给电容器充电仍能达到闪光灯击穿电压 $80\UV$，所以闪光灯仍然发光，闪光一次通过的电荷量不变，故 D 正确。
}

\item
%% number: 9
%% paperName: 2018年江苏高考第 9 题 4 分
%% typeId: 2
%% type: 多选题
%% chapter: 电磁感应
%% point: 电磁感应中的变加速
%% method: 无
%% score: 4
%% degree: 350
%% duplicateId: 0
%% body:
如图所示，竖直放置的 $\cap$ 形光滑导轨宽为 $L$，矩形匀强磁场 Ⅰ、Ⅱ 的高和间距均为 $d$，磁感应强度为 $B$。质量为 $m$ 的水平金属杆由静止释放，进入磁场 Ⅰ 和 Ⅱ 时的速度相等。金属杆在导轨间的电阻为 $R$，与导轨接触良好，其余电阻不计，重力加速度为 $g$。金属杆\xzanswer{BC}

\onepicture[w=5.2cm]{figs/09.png}

\fourchoices[answer=BC]
{刚进入磁场 Ⅰ 时加速度方向竖直向下}
{穿过磁场 Ⅰ 的时间大于在两磁场之间的运动时间}
{穿过两磁场产生的总热量为 $4mgd$}
{释放时距磁场 Ⅰ 上边界的高度 $h$ 可能小于 $\frac{m^{2}gR^{2}}{2B^{4}L^{4}}$}

%% answer: BC
%% memo:
\memoanswer{
由于金属棒进入两个磁场的速度相等，而穿出磁场后金属杆做加速度为 $g$ 的加速运动，所以金属杆进入磁场时应做减速运动，故 A 错误；对金属杆受力分析，由 $\frac{B^{2}L^{2}v}{R}-mg=ma$ 可知，金属杆做加速度减小的减速运动，其进出磁场的 $v-t$ 图象如图所示，由于 $0\sim t_{1}$ 和 $t_{1}\sim t_{2}$ 的图线与 $t$ 轴包围的面积相等（都为 $d$），所以 $t_{1}>(t_{2}-t_{1})$，故 B 正确；从进入 Ⅰ 磁场到进入 Ⅱ 磁场之前过程中，由能量守恒，金属棒减小的机械能全部转化为焦耳热，所以 $Q_{1}=mg\cdot2d$，所以穿过两个磁场过程中产生的热量为 $4mgd$，故 C 正确；若金属杆进入磁场做匀速运动，则 $\frac{B^{2}L^{2}v}{R}-mg=0$，得 $v=\frac{mgR}{B^{2}L^{2}}$，由前面分析可知金属杆进入磁场的速度大于 $\frac{mgR}{B^{2}L^{2}}$，由 $h=\frac{v^{2}}{2g}$ 得金属杆进入磁场的高度应大于 $\frac{m^{2}g^{2}R^{2}}{2gB^{4}L^{4}}=\frac{m^{2}gR^{2}}{2B^{4}L^{4}}$，故 D 错误。

\onepicture[w=4.2cm]{figs/09_an.jpg}
}

\item
%% number: 10
%% paperName: 2018年江苏高考第 10 题 8 分
%% typeId: 4
%% type: 实验题
%% chapter: 电路
%% point: 测电动势与内阻
%% method: 无
%% score: 8
%% degree: 550
%% duplicateId: 0
%% body:
一同学测量某干电池的电动势和内阻。

\begin{enumerate}
	\item
	如图所示是该同学正准备接入最后一根导线（图中虚线所示）时的实验电路。请指出图中在器材操作上存在的两个不妥之处\tkanswer[2.2cm]{开关未断开}；\tkanswer[2.2cm]{电阻箱阻值为零}。
	
	\onepicture[w=7.2cm, align=right]{figs/10a.png}
	
	\item
	实验测得的电阻箱阻值 $R$ 和电流表示数 $I$，以及计算的 $\frac{1}{I}$ 数据见下表：
	
	\begin{center}
	\begin{tblr}{|c|c|c|c|c|c|}
	\hline
	$R/\UO[a]$ & 8.0 & 7.0 & 6.0 & 5.0 & 4.0 \\
	\hline
	$I/\UA[a]$ & 0.15 & 0.17 & 0.19 & 0.22 & 0.26 \\
	\hline
	$\frac{1}{I}/\UA^{-1}$ & 6.7 & 6.0 & 5.3 & 4.5 & 3.8 \\
	\hline
	\end{tblr}
	\end{center}
	
	根据表中数据，在答题卡的方格纸上作出 $R-\frac{1}{I}$ 关系图像。
	
	\onepicture[w=6.4cm, align=right]{\includesvg{figs/10b}}
	
	由图像可计算出该干电池的电动势为\tkanswer[1.4cm]{1.4}$\UV$；内阻为\tkanswer[1.4cm]{1.2}$\UO$。
	
	\item
	为了得到更准确的测量结果，在测出上述数据后，该同学将一只量程为 $100\UmV$ 的电压表并联在电流表的两端。调节电阻箱，当电流表的示数为 $0.33\UA$ 时，电压表的指针位置如图所示，则该干电池的电动势应为\tkanswer[1.4cm]{1.4}$\UV$；内阻应为\tkanswer[1.4cm]{1.0}$\UO$。
	
	\onepicture[w=7.0cm, align=right]{figs/10c.png}
\end{enumerate}

%% answer:
\jdanswer{
\begin{enumerate}
	\item
	开关未断开，电阻箱阻值为零
	
	\item
	$1.4$（$1.30\sim1.44$ 都算对），$1.2$（$1.0\sim1.4$ 都算对）
	
	\item
	$1.4$（结果与（2）问第一个空格一致），$1.0$（结果比（2）问第二个空格小 $0.2$）
\end{enumerate}
}
%% memo:
\memoanswer{
（1）连接电路时电源应与电路断开，所以开关要断开；另一错误是电阻箱接入电路的电阻是零，这样容易烧坏电流表和电源。

（2）将数据描点连线，作出一条倾斜的直线。由闭合电路欧姆定律 $E=I(R+r)$ 得 $R=E\cdot\frac{1}{I}-r$，所以图线的斜率表示电源电动势 $E=\frac{8-(-1.2)}{6.7-0}\UV=1.37\UV$，截距绝对值表示内阻，$r=0.4\times3.0\UO=1.20\UO$；用电压表与电流表并联，可测得电流表的内阻 $R_{A}=\frac{U_{V}}{I_{A}}=\frac{66.0\times10^{-3}}{33\times10^{-3}}\UO=0.20\UO$，考虑电表内阻对实验的影响，则 $E=I(R+R_{A}+r)$，得 $R=E\cdot\frac{1}{I}-(R_{A}+r)$，所以图线的斜率仍表示电动势，电动势的准确值为 $1.37\UV$，图线的截距表示 $(R_{A}+r)$，所以内阻精确值为 $r=(1.20-0.20)\UO=1.00\UO$。

图像如图所示：

\onepicture[w=6.4cm]{\includesvg{figs/10_an}}
}

\item
%% number: 11
%% paperName: 2018年江苏高考第 11 题 10 分
%% typeId: 4
%% type: 实验题
%% chapter: 运动和力
%% point: 牛顿定律的应用
%% method: 无
%% score: 10
%% degree: 450
%% duplicateId: 0
%% body:
某同学利用如图所示的实验装置来测量重力加速度 $g$。细绳跨过固定在铁架台上的轻质滑轮，两端各悬挂一只质量为 $M$ 的重锤。实验操作如下：

\begin{steps}
	\item
	用米尺量出重锤 1 底端距地面的高度 $H$；
	\item
	在重锤 1 上加上质量为 $m$ 的小钩码；
	\item
	左手将重锤 2 压在地面上，保持系统静止。释放重锤 2，同时右手开启秒表，在重锤 1 落地时停止计时，记录下落时间；
	\item
	重复测量 3 次下落时间，取其平均值作为测量值 $t$。
\end{steps}

\onepicture[w=5.6cm, align=right]{figs/11.png}

请回答下列问题
\begin{enumerate}
	\item
	步骤④可以减小对下落时间 $t$ 测量的\tkanswer[1.5cm]{偶然}（选填“偶然”或“系统”）误差。
	\item
	实验要求小钩码的质量 $m$ 要比重锤的质量 $M$ 小很多，主要是为了
	
	\fourchoices[answer=B]
	{使 $H$ 测得更准确}
	{使重锤 1 下落的时间长一些}
	{使系统的总质量近似等于 $2M$}
	{使细绳的拉力与小钩码的重力近似相等}
	\item
	滑轮的摩擦阻力会引起实验误差。现提供一些橡皮泥用于减小该误差，可以怎么做？
	\item
	使用橡皮泥改进实验后，重新进行实验测量，并测出所用橡皮泥的质量为 $m_{0}$。用实验中的测量量和已知量表示 $g$，得 $g=$\tkanswer[3.5cm]{$\frac{2(2M+m+m_{0})H}{mt^{2}}$}。
\end{enumerate}

%% answer:
\jdanswer{
\begin{enumerate}
	\item
	偶然
	
	\item
	B
	
	\item
	在重锤 1 上粘上橡皮泥，调整橡皮泥质量直至轻拉重锤 1 能观察到其匀速下落
	
	\item
	$\frac{2(2M+m+m_{0})H}{mt^{2}}$
\end{enumerate}
}
%% memo:
\memoanswer{
（1）时间测量的误差是人为操作快慢和读数问题带来的误差，所以属于偶然误差。

（2）由于自由落体的加速度较大，下落 $H$ 高度的时间较短，为了减小测量时间的实验误差，就要使重锤下落的时间长一些，因此系统下降的加速度要小，所以小钩码的质量要比重锤的质量小很多。

（3）为了消除滑轮的摩擦阻力，可用橡皮泥粘在重锤 1 上，轻拉重锤放手后若系统做匀速运动，则表示平衡了阻力。

（4）由牛顿第二定律有 $mg=(M+m_{0}+M+m)a$，又 $H=\frac{1}{2}at^{2}$，解得 $g=\frac{2(2M+m_{0}+m)H}{mt^{2}}$。
}

\item
%% number: 12
%% paperName: 2018年江苏高考第 12 题 4 分
%% typeId: 6
%% type: 计算题
%% chapter: 运动和力
%% point: 动量定理
%% method: 无
%% score: 4
%% degree: 750
%% duplicateId: 0
%% body:
\textbf{A．}（选修模块 3-3）

\begin{enumerate}
	\item
	如图所示，一支温度计的玻璃泡外包着纱布，纱布的下端浸在水中。纱布中的水在蒸发时带走热量，使温度计示数低于周围空气温度。当空气温度不变，若一段时间后发现该温度计示数减小，则\xzanswer{A}
	
	\onepicture[w=2.4cm]{figs/12a.png}
	
	\fourchoices[answer=A]
	{空气的相对湿度减小}
	{空气中水蒸汽的压强增大}
	{空气中水的饱和气压减小}
	{空气中水的饱和气压增大}
	\item
	一定量的氧气贮存在密封容器中，在 $T_{1}$ 和 $T_{2}$ 温度下其分子速率分布的情况见右表。则 $T_{1}$\tkanswer[1.2cm]{大于}（选填“大于”“小于”或“等于”）$T_{2}$。若约 $10\%$ 的氧气从容器中泄漏，泄漏前后容器内温度均为 $T_{1}$，则在泄漏后的容器中，速率处于 $400\sim500\Ums$ 区间的氧气分子数占总分子数的百分比\tkanswer[1.2cm]{等于}（选填“大于”“小于”或“等于”）$18.6\%$。
	
	\begin{center}
	\begin{tblr}{|c|c|c|}
	\hline
	\SetCell[r=2]{c} 速率区间/$\Umdsn$ & \SetCell[c=2]{c} 各速率区间的分子数占总分子数的百分比/\% & \\
	\hline
	 & 温度 $T_{1}$ & 温度 $T_{2}$ \\
	\hline
	100 以下 & 0.7 & 1.4 \\
	\hline
	100\textasciitilde200 & 5.4 & 8.1 \\
	\hline
	200\textasciitilde300 & 11.9 & 17.0 \\
	\hline
	300\textasciitilde400 & 17.4 & 21.4 \\
	\hline
	400\textasciitilde500 & 18.6 & 20.4 \\
	\hline
	500\textasciitilde600 & 16.7 & 15.1 \\
	\hline
	600\textasciitilde700 & 12.9 & 9.2 \\
	\hline
	700\textasciitilde800 & 7.9 & 4.5 \\
	\hline
	800\textasciitilde900 & 4.6 & 2.0 \\
	\hline
	900 以上 & 3.9 & 0.9 \\
	\hline
	\end{tblr}
	\end{center}
	\item
	如图所示，一定质量的理想气体在状态 A 时压强为 $2.0\times10^{5}\UPa$，经历 A$\to$B$\to$C$\to$A 的过程，整个过程中对外界放出 $61.4\UJ$ 热量。求该气体在 A$\to$B 过程中对外界所做的功。
	
	\onepicture[w=5.0cm, align=right]{figs/12b.png}
\end{enumerate}

\textbf{B．}（选修模块 3-4）

\begin{enumerate}
	\item
	梳子在梳头后带上电荷，摇动这把梳子在空中产生电磁波。该电磁波\xzanswer{AD}
	
	\fourchoices[answer=AD]
	{是横波}
	{不能在真空中传播}
	{只能沿着梳子摇动的方向传播}
	{在空气中的传播速度约为 $3\times10^{8}\Ums$}
	\item
	两束单色光 A、B 的波长分别为 $\lambda_{A}$、$\lambda_{B}$，且 $\lambda_{A}>\lambda_{B}$，则\tkanswer[0.8cm]{A}（选填“A”或“B”）在水中发生全反射时的临界角较大。用同一装置进行杨氏双缝干涉实验时，可以观察到\tkanswer[0.8cm]{A}（选填“A”或“B”）产生的条纹间距较大。
	\item
	一列简谐横波沿 $x$ 轴正方向传播，在 $x=0$ 和 $x=0.6\Um$ 处的两个质点 A、B 的振动图象如图所示。已知该波的波长大于 $0.6\Um$，求其波速和波长。
	
	\onepicture[w=7.2cm, align=right]{figs/12c.png}
\end{enumerate}

\textbf{C．}（选修模块 3-5）

\begin{enumerate}
	\item
	已知 A 和 B 两种放射性元素的半衰期分别为 $T$ 和 $2T$，则相同质量的 A 和 B 经过 $2T$ 后，剩有的 A 和 B 质量之比为\xzanswer{B}
	
	\fourchoices[answer=B]
	{$1:4$}
	{$1:2$}
	{$2:1$}
	{$4:1$}
	\item
	光电效应实验中，用波长为 $\lambda_{0}$ 的单色光 A 照射某金属板时，刚好有光电子从金属表面逸出。当波长为 $\frac{\lambda_{0}}{2}$ 的单色光 B 照射该金属板时，光电子的最大初动能为\tkanswer[2.0cm]{$\frac{hc}{\lambda_{0}}$}，A、B 两种光子的动量之比为\tkanswer[1.2cm]{$1:2$}。（已知普朗克常量为 $h$、光速为 $c$）
	\item
	如图所示，悬挂于竖直弹簧下端的小球质量为 $m$，运动速度的大小为 $v$，方向向下。经过时间 $t$，小球的速度大小为 $v$，方向变为向上。忽略空气阻力，重力加速度为 $g$，求该运动过程中，小球所受弹簧弹力冲量的大小。
	
	\onepicture[w=2.6cm, align=right]{figs/12d.png}
\end{enumerate}

%% answer:
\jdanswer{
\begin{enumerate}
	\item
	A：（1）A；（2）大于，等于；（3）气体对外界做的功为 $138.6\UJ$。
	\item
	B：（1）AD；（2）A，A；（3）$v=2\Ums$；$\lambda=0.8\Um$。
	\item
	C：（1）B；（2）$\frac{hc}{\lambda_{0}}$，$1:2$；（3）$2mv+mgt$。
\end{enumerate}
}
%% memo:
\memoanswer{
\textbf{A．}（1）温度计示数减小说明蒸发加快，空气中水蒸气的压强减小，故 B 错误；因空气中水的饱和汽压只与温度有关，空气温度不变，所以水的饱和汽压不变，故 C、D 错误；由相对湿度的定义，空气的相对湿度减小，故 A 正确。（2）分子速率分布与温度有关，温度升高，分子的平均速率增大，速率大的分子数所占比例增加，速率小的分子数所占比例减小，所以 $T_{1}$ 大于 $T_{2}$；泄漏前后容器内温度不变，则在泄漏后的容器中，速率处于 $400\sim500\Ums$ 区间的氧气分子数占总分子数的百分比不变，仍为 $18.6\%$。（3）整个过程中，外界对气体做的功为 $W=W_{AB}+W_{CA}$，且 $W_{CA}=p_{A}(V_{C}-V_{A})$。由热力学第一定律 $\Delta U=Q+W$ 得 $W_{AB}=-(Q+W_{CA})$。代入数据得 $W_{AB}=-138.6\UJ$，即气体对外界做的功为 $138.6\UJ$。

\textbf{B．}（1）摇动的梳子在空中产生电磁波，电磁波是横波，故 A 正确；电磁波能在真空中传播，故 B 错误；电磁波传播的方向与振动方向垂直，故 C 错误；电磁波在空气中传播的速度约为光速，故 D 正确。（2）波长越长，频率越小，折射率越小，由临界角公式 $\sin C=\frac{1}{n}$，可知波长越大临界角越大，所以 A 光的临界角大；双缝干涉条纹的间距 $\Delta x=\frac{l}{d}\lambda$，因为 A 光的波长较长，所以 A 光产生的条纹间距较大。（3）由图象可知，周期 $T=0.4\Us$，由于波长大于 $0.6\Um$，由图象可知，波从 A 到 B 的传播时间 $\Delta t=0.3\Us$，波速 $v=\frac{\Delta x}{\Delta t}$，代入数据得 $v=2\Ums$，波长 $\lambda=vT$，代入数据得 $\lambda=0.8\Um$。

\textbf{C．}（1）A 的半衰期为 $T$，经过 $2T$ 后剩余质量为原来的 $\frac{1}{4}$；B 的半衰期为 $2T$，经过 $2T$ 后剩余质量为原来的 $\frac{1}{2}$，故剩余质量之比为 $\frac{1}{4}:\frac{1}{2}=1:2$，故 B 正确。（2）根据光电效应方程 $E_{\text{k}}=h\nu-W_{0}$，又 $\nu=\frac{c}{\lambda}$，所以有 $0=\frac{hc}{\lambda_{0}}-W_{0}$，$E_{\text{k}}=\frac{2hc}{\lambda_{0}}-W_{0}$，解得 $E_{\text{k}}=\frac{hc}{\lambda_{0}}$；又光子动量 $p=\frac{h}{\lambda}$，所以 A、B 两种光子的动量之比为 $1:2$。（3）取向上为正方向，由动量定理得 $mv-(-mv)=I$，且 $I=(\overline{F}-mg)t$，解得 $I_{F}=\overline{F}t=2mv+mgt$。
}

\item
%% number: 13
%% paperName: 2018年江苏高考第 13 题 15 分
%% typeId: 1
%% type: 单选题
%% chapter: 磁场
%% point: 安培力
%% method: 无
%% score: 15
%% degree: 650
%% duplicateId: 0
%% body:
如图所示，两条平行的光滑金属导轨所在平面与水平面的夹角为 $\theta$，间距为 $d$。导轨处于匀强磁场中，磁感应强度大小为 $B$，方向与导轨平面垂直。质量为 $m$ 的金属棒被固定在导轨上，距底端的距离为 $s$，导轨与外接电源相连，使金属棒通有电流。金属棒被松开后，以加速度 $a$ 沿导轨匀加速下滑，金属棒中的电流始终保持恒定，重力加速度为 $g$。求下滑到底端的过程中，金属棒

\begin{enumerate}
	\item
	末速度的大小 $v$；
	\item
	通过的电流大小 $I$；
	\item
	通过的电荷量 $Q$。
\end{enumerate}

\onepicture[w=6.4cm, align=right]{figs/13.png}

%% answer: IBB
%% memo:
\memoanswer{
（1）匀加速直线运动 $v^{2}=2as$，解得 $v=\sqrt{2as}$。

（2）安培力 $F_{\text{安}}=IdB$，金属棒所受合力 $F=mg\sin\theta-F_{\text{安}}$，牛顿运动定律 $F=ma$，解得 $I=\frac{mg(\sin\theta-a)}{dB}$。

（3）运动时间 $t=\frac{v}{a}$，电荷量 $Q=It$，解得 $Q=\frac{\sqrt{2as}\,m(g\sin\theta-a)}{dBa}$。
}

\item
%% number: 14
%% paperName: 2018年江苏高考第 14 题 16 分
%% typeId: 6
%% type: 计算题
%% chapter: 功和能
%% point: 整体机械能守恒
%% method: 无
%% score: 16
%% degree: 450
%% duplicateId: 0
%% body:
如图所示，钉子 A、B 相距 $5l$，处于同一高度。细线的一端系有质量为 $M$ 的小物块，另一端绕过 A 固定于 B。质量为 $m$ 的小球固定在细线上 C 点，B、C 间的线长为 $3l$。用手竖直向下拉住小球，使小球和物块都静止，此时 BC 与水平方向的夹角为 $53^{\circ}$。松手后，小球运动到与 A、B 相同高度时的速度恰好为零，然后向下运动。忽略一切摩擦，重力加速度为 $g$，取 $\sin53^{\circ}=0.8$，$\cos53^{\circ}=0.6$。求：

\begin{enumerate}
	\item
	小球受到手的拉力大小 $F$；
	\item
	物块和小球的质量之比 $M:m$；
	\item
	小球向下运动到最低点时，物块 $M$ 所受的拉力大小 $T$。
\end{enumerate}

\onepicture[w=5.6cm, align=right]{\includesvg{figs/14}}

%% answer:
\jdanswer{
\begin{enumerate}
	\item
	$F=\frac{5}{3}Mg-mg$
	\item
	$\frac{M}{m}=\frac{6}{5}$
	\item
	$T=\frac{8mMg}{5(m+M)}$（$T=\frac{48}{55}mg$ 或 $T=\frac{8}{11}Mg$）
\end{enumerate}
}
%% memo:
\memoanswer{
（1）设小球受 AC、BC 的拉力分别为 $F_{1}$、$F_{2}$，由平衡条件得 $F_{1}\sin53^{\circ}=F_{2}\cos53^{\circ}$，$F+mg=F_{1}\cos53^{\circ}+F_{2}\sin53^{\circ}$，且 $F_{1}=Mg$，解得 $F=\frac{5}{3}Mg-mg$。

（2）小球运动到与 A、B 相同高度过程中，小球上升高度 $h_{1}=3l\sin53^{\circ}$，物块下降高度 $h_{2}=2l$，由机械能守恒定律得 $mgh_{1}=Mgh_{2}$，解得 $\frac{M}{m}=\frac{6}{5}$。

（3）根据机械能守恒定律，小球回到起始点。设此时 AC 方向的加速度大小为 $a$，重物受到的拉力为 $T$，牛顿运动定律 $Mg-T=Ma$，小球受 AC 的拉力 $T'=T$，牛顿运动定律 $T'-mg\cos53^{\circ}=ma$，解得 $T=\frac{8mMg}{5(m+M)}$（$T=\frac{48}{55}mg$ 或 $T=\frac{8}{11}Mg$）。
}

\item
%% number: 15
%% paperName: 2018年江苏高考第 15 题 16 分
%% typeId: 6
%% type: 计算题
%% chapter: 磁场
%% point: 带电粒子在磁场中的运动
%% method: 无
%% score: 16
%% degree: 350
%% duplicateId: 0
%% body:
如图所示，真空中四个相同的矩形匀强磁场区域，高为 $4d$，宽为 $d$，中间两个磁场区域间隔为 $2d$，中轴线与磁场区域两侧相交于 O、O$'$ 点，各区域磁感应强度大小相等。某粒子质量为 $m$、电荷量为 $+q$，从 O 沿轴线射入磁场。当入射速度为 $v_{0}$ 时，粒子从 O 上方 $\frac{d}{2}$ 处射出磁场。取 $\sin53^{\circ}=0.8$，$\cos53^{\circ}=0.6$。

\begin{enumerate}
	\item
	求磁感应强度大小 $B$；
	\item
	入射速度为 $5v_{0}$ 时，求粒子从 O 运动到 O$'$ 的时间 $t$；
	\item
	入射速度仍为 $5v_{0}$，通过沿轴线 OO$'$ 平移中间两个磁场（磁场不重叠），可使粒子从 O 运动到 O$'$ 的时间增加 $\Delta t$，求 $\Delta t$ 的最大值。
\end{enumerate}

\onepicture[w=5.4cm, align=right]{figs/15.png}

%% answer:
\jdanswer{
\begin{enumerate}
	\item
	$B=\frac{4mv_{0}}{qd}$
	\item
	$t=\left(\frac{53\pi+72}{180}\right)\frac{d}{v_{0}}$
	\item
	$\Delta t_{m}=\frac{d}{5v_{0}}$
\end{enumerate}
}
%% memo:
\memoanswer{
（1）粒子在磁场中，由洛伦兹力提供向心力得 $qvB=m\frac{v^{2}}{r}$，解得 $r_{0}=\frac{mv_{0}}{qB}$，由几何关系知 $r_{0}=\frac{d}{4}$，联立解得 $B=\frac{4mv_{0}}{qd}$。

（2）设粒子在矩形磁场中的偏转角为 $\alpha$，由 $d=r\sin\alpha$，得 $\sin\alpha=\frac{4}{5}$，即 $\alpha=53^{\circ}$。在一个矩形磁场中的运动时间 $t_{1}=\frac{\alpha}{360^{\circ}}\cdot\frac{2\pi m}{qB}$，解得 $t_{1}=\frac{53\pi d}{720v_{0}}$；直线运动的时间 $t_{2}=\frac{2d}{v}$，解得 $t_{2}=\frac{2d}{5v_{0}}$；则 $t=4t_{1}+t_{2}=\left(\frac{53\pi+72}{180}\right)\frac{d}{v_{0}}$。

（3）将中间两磁场分别向中央移动距离 $x$，粒子向上的偏移量 $y=2r(1-\cos\alpha)+x\tan\alpha$，由 $y\leq2d$，解得 $x\leq\frac{3}{4}d$，则当 $x_{m}=\frac{3}{4}d$ 时，$\Delta t$ 有最大值。粒子直线运动路程的最大值 $s_{m}=\frac{2x_{m}}{\cos\alpha}+(2d-2x_{m})=3d$，增加路程的最大值 $\Delta s_{m}=s_{m}-2d=d$，增加时间的最大值 $\Delta t_{m}=\frac{\Delta s_{m}}{v}=\frac{d}{5v_{0}}$。
}

\end{enumerate}

\end{document}
